% English labels and lesson macros.
\renewcommand{\contentsname}{Contents}
\renewcommand{\figurename}{Figure}
\renewcommand{\tablename}{Table}
\renewcommand{\chaptername}{Chapter}
\renewcommand{\appendixname}{Appendix}
\hypersetup{pdftitle={Understanding simpleKafka from Scratch},pdfauthor={simpleKafka teaching project}}

\newenvironment{contract}{\begin{quote}\small\noindent\textbf{Behavioral contract and invariants.}\ }{\end{quote}}
\newcommand{\lessoncommand}[1]{\par\noindent\textbf{Commands and progression.} Use \coursecmd{./gradlew :stepTest -Pstep=#1} to diagnose this step; \coursecmd{./gradlew :test -Pstep=#1} runs cumulatively from Step 1 through this step. A chapter is complete only when the cumulative tests through its final step (4, 8, 12, 16, 20, 24, 28, or 31) all pass.}
\newcommand{\lessonquestions}[2]{\subsubsection*{Check your understanding}\begin{enumerate}\item #1\item #2\end{enumerate}}
\newcommand{\lessonhints}[3]{\subsubsection*{Tiered hints (not the answer)}\begin{hintlevels}\item[Level 1] #1\item[Level 2] #2\item[Level 3] #3\end{hintlevels}}
\newcommand{\lessonsection}[2]{\section{Step #1: #2}\label{step:#1}}
\newcommand{\stepref}[2]{\hyperref[step:#1]{Step #1 (#2)}}
\newcommand{\expected}[1]{\begin{quote}\small\textbf{Observable examples and test points.} #1\end{quote}}
\newcommand{\providededit}[2]{\begin{quote}\small\textbf{Starting point and edit location.}\par Already provided: #1\par Implement in: #2\end{quote}}
\newcommand{\pitfalls}[1]{\subsubsection*{Common pitfalls and differences from real Kafka}#1}
\newcommand{\goals}[2]{\subsubsection*{Goals and prerequisites}Goal: #1\par Before you begin, be able to explain: #2}
\newcommand{\background}[1]{\subsubsection*{Background and reasoning}#1}
\newcommand{\lessoncontract}[1]{\begin{contract}#1\end{contract}}
